\section{Outer rows and inherited fibres}\label{sec:outer}

The small-radical input is Theorem~3 of \cite{Paper1}.
Appendix~\ref{app:sources} states it, and Appendix~\ref{app:lowrad}
records the shell substitution. With \(H=Y/2\) and
\(V=Y^{5/3-\epsilon_0}\), it gives
\[
\#\{n:Y/2<n\le Y,\ \operatorname{rad}(n(n+1))\le
Y^{5/3-\epsilon_0}\}
\ll_\eta Y^{2/3-S_0+\eta}.
\AutoTag{eq:old:4_2}
\]
The threshold in that application is large enough that
\(Y\ge4\) and \(Y^{1/3+\epsilon_0}>4\). These are exactly the required
conditions \(H\ge2\) and \(V<H^2\).

Let \(Q_{\rm sf}\) be a positive squarefree integer. For a positive
integer \(Q\), let \(\omega(Q)\) denote the number of distinct prime
divisors of \(Q\). Apply the support-fibre estimate
\cite[Lemma~2.1]{Paper1} with \(X=n\):
\[
\#\{1\le m<n:\operatorname{rad}(m(m+1))\mid Q_{\rm sf}\}
 \ll_\eta Y^{2\eta}\log(2Y).
\AutoTag{eq:old:4_3}
\]
Here \(Q_{\rm sf}=\operatorname{rad}(n(n+1))\) is squarefree and at most
\(Y(Y+1)\). The source equality of radicals gives the divisibility
condition for every partner. The bound follows from
\(2^{\omega(Q_{\rm sf})}\ll_\eta Q_{\rm sf}^\eta\). Choose the
internal exponents smaller than the requested final \(\varepsilon\).
Absorbing the logarithm then proves the \(O_R\) row bound.

\begin{proposition}[Five outer rows]\label{prop:outerrows}
For every \(\varepsilon>0\),
\[
\begin{aligned}
\#O_R,\ \#O_M,\ \#O_{\operatorname{rad}M},\ \#Q1
  &\ll_\varepsilon Y^{2/3-S_0+\varepsilon},\\
\#Q0&\ll_\varepsilon Y^\varepsilon.
\end{aligned}\AutoTag{eq:old:4_4}
\]
\end{proposition}

\begin{proof}
The \(O_R\) estimate is \eqref{eq:old:4_2} followed by the partner-fibre bound \eqref{eq:old:4_3}.
We apply the latter with \(X=n\) and with the squarefree common radical
as \(Q\). The defining equality of radicals implies the required
divisibility \(\operatorname{rad}(m(m+1))\mid Q\). Each logarithm and divisor factor uses a
separate portion of the arbitrary \(\varepsilon\).

For \(O_M\), fix \(M\le Y^{2/3-\epsilon_0}\). Write
\(P=ab\), \(Q_{\rm prod}=cd\), so \(PQ_{\rm prod}=M\),
\(P<Q_{\rm prod}\), and \(Q_{\rm prod}-P=L>0\).
There are at most \(d_4(M)\ll_\eta M^\eta\) ordered choices of
\((a,b,c,d)\), where
\[
d_4(M)=\#\{(a,b,c,d)\in\mathbb Z_{>0}^4:abcd=M\}.
\]
For each, the source identity
\[
Q_{\rm prod}m(m+1)=Pn(n+1)
\]
is equivalent, after setting \(x=2n+1\) and \(y=2m+1\), to
\[
P x^2+(Q_{\rm prod}-P)=Q_{\rm prod}y^2.
\AutoTag{eq:old:4_5}
\]
Appendix~\ref{app:fibre} bounds the number of such pairs by \(Y^\eta\).
The coefficients and nonzero shift have height bounded by a fixed power
of \(Y\). For the theorem's variable bound take \(x_0=2Y+1\), since
\(2n+1\le x_0\). The source identity gives
\(M\le Y^2(Y+1)^2\le4Y^4\), so
\(P,Q_{\rm prod},|Q_{\rm prod}-P|=L\le4Y^4\le x_0^4\), and the
fixed height exponent is \(A_{\rm ht}=4\).
Thus the implied constant is uniform in the moving labels. Summing over
\(M\) proves the \(O_M\) bound.

For \(O_{\operatorname{rad}M}\), \eqref{eq:old:2_7} gives \(M\le4Y^4\). If
\(R_0=Y^{2/3-k_0}\), then
\[
\#\{M\le4Y^4:\operatorname{rad}(M)\le R_0\}
\le R_0(4Y^4)^\sigma
 \sum_{M\ge1}\frac1{\operatorname{rad}(M)M^\sigma}.
\AutoTag{eq:old:4_6}
\]
For fixed \(\sigma>0\), the series is the convergent product
\[
\prod_q\left(1+\frac1{q(q^\sigma-1)}\right).
\]
Fix the final target loss. Then choose a fixed \(\sigma>0\) so small
that \(4\sigma\) is below its allocated part. Choose the divisor and
fibre exponents independently, each within its own remaining allocation.
The resulting label count is
\(O_\sigma(Y^{2/3-k_0+4\sigma})\).
For each \(M\), sum the at most \(d_4(M)\) residual factorizations and
apply the same fixed-height fibre \eqref{eq:old:4_5}. Since \(k_0=S_0\), this gives
the asserted exponent.

On \(Q0\), failure of the first three rows and the \(Q0\) predicate give
\[
Y^{q_{\max}}<M<\kappa_Y^2Y^{q_{\max}}.
\]
The interval length is at most
\((\kappa_Y^2-1)Y^{q_{\max}}\le3Y^{q_{\max}-1}<1\) for \(Y\ge256\).
Thus at most one integer \(M\) occurs. Its complete ordered
factorization sum and the fibre \eqref{eq:old:4_5} cost \(Y^\varepsilon\). We do
not declare the row empty.

On \(Q1\), enumerate \(P=ab\) and \(L=cd-ab\), and put
\(Q_{\rm prod}=cd=P+L\). Since \(M=PQ_{\rm prod}\ge P^2\),
\eqref{eq:old:2_7} gives \(P<\kappa_YY^{q_{\max}/2}\ll Y^{1/3+S_0}\).
The \(Q1\) predicate and the positivity of the integer \(L\) give
\(1\le L<Y^{1/3-l_0}=Y^{1/3-2S_0}\). Hence the number of \((P,L)\)
labels is
\[
\ll Y^{1/3+S_0}Y^{1/3-2S_0}=Y^{2/3-S_0}.
\]
Each label fixes \(Q_{\rm prod}=P+L\). Splitting \(P,Q_{\rm prod}\)
into their ordered factors costs at most
\(\tauDiv(P)\tauDiv(Q_{\rm prod})\ll_\eta Y^\eta\). Then the fixed-height fibre
\eqref{eq:old:4_5} is used once. This row carries no additional divisor sum
over \(M\).
\end{proof}

\begin{remark}
The source and theorem hypotheses in this section hold above fixed
thresholds. Their applications are recorded in
Appendix~\ref{app:lowrad}.
The finitely many remaining \(Y\) are absorbed into the implied constants
in Proposition~\ref{prop:outerrows}.
\end{remark}
